\begin{thebibliography}{99}
\bibitem{source1} Prashanth Amireddy, Ankit Garg, Neeraj Kayal, Chandan Saha, and Bhargav Thankey. Low- depth arithmetic circuit lower bounds: Bypassing set-multilinearization. In Proceedings of the 50th International Colloquium on Automata, Languages and Programming (ICALP 2023), volume 261 of Leibniz International Proceedings in Informatics (LIPIcs), pages 12:1–12:20, 2023. Full version in the Electronic Colloquium on Computational Complexity (ECCC), Technical Report TR22-151. doi:10.4230/LIPICS.ICALP.2023.12.
\bibitem{source2} Robert Andrews. Algebraic hardness versus randomness in low characteristic. In Proceedings of the 35th Annual Computational Complexity Conference (CCC 2020), volume 169 of Leibniz International Proceedings in Informatics (LIPIcs), pages 37:1–37:32, 2020. Full version at arXiv:2005.10885. doi:10.4230/LIPICS.CCC.2020.37.
\bibitem{source3} C. S. Bhargav, Sagnik Dutta, and Nitin Saxena. Improved lower bound, and proof barrier, for constant depth algebraic circuits. In Proceedings of the 47th Internationl Symposium on the Mathematical Foundations of Computer Science (MFCS 2022), volume 241 of Leibniz International Proceedings in Informatics (LIPIcs), pages 18:1–18:16, 2022. doi:10.4230/ LIPICS.MFCS.2022.18.
\bibitem{source4} Radu Curticapean, Nutan Limaye, and Srikanth Srinivasan. On the VNP-hardness of some monomial symmetric polynomials. In 42nd International Conference on Foundations of Software Technology and Theoretical Computer Science (FSTTCS 2022), volume 250 of Leibniz International Proceedings in Informatics (LIPIcs), pages 16:1–16:14, 2022. Full version in the Electronic Colloquium on Computational Complexity (ECCC), Technical Report TR22-139. doi:10.4230/LIPICS.FSTTCS.2022.16.
\bibitem{source5} Ismor Fischer. Sums of like powers of multivariate linear forms. Mathematics Magazine, 67(1):59–61, 1994. URL: http://www.jstor.org/stable/2690560.
\bibitem{source6} Philippe Flajolet and Robert Sedgewick. Analytic combinatorics. Cambridge University Press, 2009. doi:10.1017/CBO9780511801655.
\bibitem{source7} Michael A. Forbes, Amir Shpilka, and Ben Lee Volk. Succinct hitting sets and barriers to proving algebraic circuits lower bounds. In Proceedings of the 49th Annual ACM Symposium on Theory of Computing (STOC 2017), pages 653–664, 2017. Full version at arXiv:1701.05328. doi:10.1145/3055399.3055496.
\bibitem{source8} Hervé Fournier, Nutan Limaye, Srikanth Srinivasan, and Sébastien Tavenas. On the power of homogeneous algebraic formulas. Electronic Colloquium on Computational Complexity (ECCC), TR23-191, 2023. URL: https://eccc.weizmann.ac.il/report/2023/191.
\bibitem{source9} Nashlen Govindasamy, Tuomas Hakoniemi, and Iddo Tzameret. Simple hard instances for low-depth algebraic proofs. In Preliminary version in the 63rd Annual IEEE Symposium on Foundations of Computer Science (FOCS 2022), pages 188–199, 2022. Full version at arXiv:2205.07175. doi:10.1109/FOCS54457.2022.00025.
\bibitem{source10} Joshua A. Grochow, Mrinal Kumar, Michael E. Saks, and Shubhangi Saraf. Towards an algebraic natural proofs barrier via polynomial identity testing. arXiv, 1701.01717, 2017. URL: http://arxiv.org/abs/1701.01717.
\bibitem{source11} Joshua A. Grochow and Toniann Pitassi. Circuit complexity, proof complexity, and polynomial identity testing. In Proceedings of the 55th Annual IEEE Symposium on Foundations of Computer Science (FOCS 2014), pages 110–119, 2014. Full version at arXiv:1404.3820. doi:10.1109/FOCS.2014.20.
\bibitem{source12} Alan Guo, Swastik Kopparty, and Madhu Sudan. New affine-invariant codes from lifting. In Proceedings of Innovations in Theoretical Computer Science (ITCS 2013), pages 529–540, 2013. Full version at arXiv:1208.5413. doi:10.1145/2422436.2422494.
\bibitem{source13} Ankit Gupta, Pritish Kamath, Neeraj Kayal, and Ramprasad Saptharishi. Arithmetic circuits: A chasm at depth three. In Proceedings of the 54th Annual IEEE Symposium on Foundations of Computer Science (FOCS 2013), pages 578–587, 2013. Full version in the Electronic Colloquium on Computational Complexity (ECCC), Technical Report TR13-026. doi:10. 1109/FOCS.2013.68.
\bibitem{source14} Erich L. Kaltofen. Factorization of polynomials given by straight-line programs. In Silvio Micali, editor, Randomness and Computation, volume 5 of Advances in Computing Research, pages 375–412. JAI Press, Inc., Greenwich, CT, USA, 1989. URL: http://www.math.ncsu. edu/~kaltofen/bibliography/89/Ka89\_slpfac.pdf.
\bibitem{source15} Nutan Limaye, Srikanth Srinivasan, and Sébastien Tavenas. Superpolynomial lower bounds against low-depth algebraic circuits. In Preliminary version in the 62nd Annual IEEE Sym- posium on Foundations of Computer Science (FOCS 2021), pages 804–814, 2022. Full version in the Electronic Colloquium on Computational Complexity (ECCC), Technical Report TR21-081. doi:10.1109/FOCS52979.2021.00083.
\bibitem{source16} Nutan Limaye, Srikanth Srinivasan, and Sébastien Tavenas. Guest column: Lower bounds against constant-depth algebraic circuits. SIGACT News, 53(2):40–62, 2022. doi:10.1145/ 3544979.3544989.
\bibitem{source17} Henryk Minc. Evaluation of permanents. Proc. Edinburgh Math. Soc. (2), 22(1):27–32, 1979. doi:10.1017/S0013091500027760.
\bibitem{source18} Ran Raz. Elusive functions and lower bounds for arithmetic circuits. Theory of Computing, 6(1):135–177, 2010. 40th Annual ACM Symposium on Theory of Computing (STOC 2008). doi:10.4086/TOC.2010.V006A007.
\bibitem{source19} Alexander A. Razborov. Lower bounds on the size of bounded depth circuits over a complete basis with logical addition. Matematicheskie Zametki, 41(4):598–607, April 1987.
\bibitem{source20} Amir Shpilka and Avi Wigderson. Depth-3 arithmetic circuits over fields of characteristic zero. Computational Complexity, 10(1):1–27, 2001. Preliminary version in the 14th Annual IEEE Conference on Computational Complexity (CCC 1999). doi:10.1007/PL00001609.
\bibitem{source21} Roman Smolensky. Algebraic methods in the theory of lower bounds for boolean circuit complexity. In Proceedings of the 19th Annual ACM Symposium on Theory of Computing (STOC 1987), pages 77–82, 1987. doi:10.1145/28395.28404.
\bibitem{source22} Ryan Williams. Finding paths of length k in O∗ (2k ) time. Inf. Process. Lett., 109(6):315–318, 2009. doi:10.1016/J.IPL.2008.11.004.
\end{thebibliography}
